\section{Command reference}
\label{sec:command-reference}

The rest of this manual explains what each area of the CLI is for and how
the pieces fit together; this section is the complete index of every
\code{keyquorum} and \code{keyquorum-device} command and subcommand, so
this manual alone --- without a separate run of \code{--help} --- says
what the tool can do. Each row's description is the same one \code{--help}
prints for that command. Flags are not reproduced here in full; run
\code{keyquorum <command> --help} (or \code{keyquorum <command>
<subcommand> --help}) for a command's complete flag list, defaults, and
which flags conflict with or require others.

\begin{noteBox}
\code{keyquorum host} is a hidden subcommand, compiled only with
\flag{--features provider}, and is intentionally omitted here: it is a
build capability for operating a mailbox relay, not a customer-facing
command, and is never documented outside the crate's own source.
\end{noteBox}

\subsection{Identity and hardware keys}

\begin{center}
\begin{tabular}{@{}p{0.24\linewidth}p{0.70\linewidth}@{}}
\toprule
\textbf{Command} & \textbf{Purpose} \\
\midrule
\code{generate} & Generate an encryption or signing keypair. The private
  key prints to stdout once and is never written to disk by this tool;
  \flag{--register} also registers the new public key. \\
\code{register} & Register an existing public key under a label. \\
\code{list} & List registered hardware keys and live split trees. \\
\code{revoke} & Revoke a registered key: drops its pairings and
  private-bridge membership everywhere; \flag{--evict} also PSS-refreshes
  the surviving siblings of its leaf. \\
\code{remove} & Remove a registered hardware key. \\
\bottomrule
\end{tabular}
\end{center}

\subsection{Split-key trees}

\begin{center}
\begin{tabular}{@{}p{0.24\linewidth}p{0.70\linewidth}@{}}
\toprule
\textbf{Command} & \textbf{Purpose} \\
\midrule
\code{split} & Split a secret into a live tree. \flag{--leaf} builds the
  tree directly (siblings are bound automatically); \flag{--tree-spec} is
  for restoring a nested snapshot. \\
\code{bind} & Pair two nodes (establish \code{node <-> peer}), or reseal a
  leaf onto a new public key. \\
\code{add} & Insert a leaf under a parent split and reshare that parent. \\
\code{tree} & Print every live tree, one tree by id, the lowest common
  ancestor of \flag{--node} labels, or write the live spec as JSON. \\
\code{tree publish} & Upload a tree's public topology (no sealed shares) to
  the relay. \\
\code{tree restructure} & Announce a restructured tree: one authenticated
  envelope per active leaf at the next public generation. Unlike
  \code{publish}, recipients can verify who authorized the change. \\
\code{tree countersign} & Countersign a pending restructure proposal (key
  file, or \flag{--device}/\flag{--slot}). \\
\code{tree fetch} & Download the slice a pull key is allowed to see and
  merge it into this store. \\
\code{reconstruct} & Reconstruct a key's secret from raw shares (key
  files, container slots, or an interactive prompt). \\
\bottomrule
\end{tabular}
\end{center}

\subsection{Password vault}

\begin{center}
\begin{tabular}{@{}p{0.24\linewidth}p{0.70\linewidth}@{}}
\toprule
\textbf{Command} & \textbf{Purpose} \\
\midrule
\code{vault add} & Store a new credential, optionally PIN-protected. \\
\code{vault get} & Retrieve a stored credential. \\
\bottomrule
\end{tabular}
\end{center}

\subsection{Bridges}

\begin{center}
\begin{tabular}{@{}p{0.24\linewidth}p{0.70\linewidth}@{}}
\toprule
\textbf{Command} & \textbf{Purpose} \\
\midrule
\code{bridge allow} & Grant \flag{--node} permission to form a cross-branch
  link with \flag{--peer}. \\
\code{bridge deny} & Revoke that permission and drop any established link
  between them. \\
\code{bridge add} & Establish a pairing, if either node's whitelist allows
  it. \\
\code{bridge remove} & Tear down an established pairing; the whitelist
  entry is left intact. \\
\code{bridge list} & List whitelist entries and established pairings. \\
\code{bridge private create} & Create a private sign bridge: writes every
  recipient's \file{.kqpb} envelope first, then commits this store, so a
  failure at any point can simply be retried. \\
\code{bridge private list} & List private bridges, optionally for one
  split tree. \\
\code{bridge private show} & Show one private bridge by uid. \\
\code{bridge private events} & Print bridge-change events. \\
\code{bridge private import} & Import a per-person bridge package into this
  store. \\
\code{bridge private remove-member} & Drop a signing member: writes
  rotation/destroy packages first, then commits. \\
\bottomrule
\end{tabular}
\end{center}

\subsection{Authenticated updates}

\begin{center}
\begin{tabular}{@{}p{0.24\linewidth}p{0.70\linewidth}@{}}
\toprule
\textbf{Command} & \textbf{Purpose} \\
\midrule
\code{reissue} & Announce a replaced hardware key: writes one authenticated
  envelope per store that holds the old key, so every recipient converges
  on the same token. \\
\code{updates} & Print the authenticated organization updates this store
  has applied. \\
\bottomrule
\end{tabular}
\end{center}

\subsection{Files and credentials}

\begin{center}
\begin{tabular}{@{}p{0.24\linewidth}p{0.70\linewidth}@{}}
\toprule
\textbf{Command} & \textbf{Purpose} \\
\midrule
\code{access password} & Lock (\flag{--state 0}) or unlock (\flag{--state
  1}) a password-protected file, optionally with a PIN and an expiry. \\
\code{access quorum} & Lock (\flag{--state 0}) or unlock (\flag{--state
  1}) a hardware-key-quorum-protected file, or print its status
  (\flag{--status}) --- custody mode, minimum physical devices, unlock
  approval, and requirement tree. \\
\code{export credential} & Export a vault credential, sealed to a
  recipient's public key. \\
\code{export file} & Export a password-locked file, sealed to a
  recipient's public key. \\
\code{share create-credential} & Create a time-limited, revocable share
  link for a vault credential. \\
\code{share create-file} & Create a time-limited, revocable share link for
  a password-locked file. \\
\code{share redeem-credential} & Redeem a credential share token (prompted
  interactively, never as an argument). \\
\code{share redeem-file} & Redeem a file share token (prompted
  interactively, never as an argument). \\
\code{share revoke-credential} & Revoke a credential share. \\
\code{share revoke-file} & Revoke a file share. \\
\code{pin relock} & End a cached one-time PIN unlock window immediately. \\
\bottomrule
\end{tabular}
\end{center}

\subsection{Signing}

\begin{center}
\begin{tabular}{@{}p{0.24\linewidth}p{0.70\linewidth}@{}}
\toprule
\textbf{Command} & \textbf{Purpose} \\
\midrule
\code{sign} & Sign a message with a private bridge plus a personal signing
  key. \\
\code{verify} & Verify an Ed25519 signature, standalone or against a
  private bridge's roster (\flag{--bridge-uid}). \\
\bottomrule
\end{tabular}
\end{center}

\subsection{Delivery and the mailbox relay}

\begin{center}
\begin{tabular}{@{}p{0.24\linewidth}p{0.70\linewidth}@{}}
\toprule
\textbf{Command} & \textbf{Purpose} \\
\midrule
\code{deliver send} & Seal a file to a registered label, signed with your
  signing key. \\
\code{deliver open} & Open a letter addressed to you: verify the sender,
  keep (or \flag{--reject}) the file, and seal a signed acknowledgement
  back. \\
\code{deliver ack} & Check an acknowledgement sealed back to you. \\
\code{relay push} & Upload every \file{.kqpb} in a directory to the relay;
  also replaces public-tree documents from trees stored in \flag{--db}. \\
\code{relay pull} & Download envelopes and the public-tree slice for a
  pull key; \flag{--import} installs each envelope into \flag{--db}
  directly. \\
\code{loadkey} & Check a relay API key with \code{POST /keycheck} and
  store it (and its hash) on this instance for later relay commands to
  reuse. \\
\bottomrule
\end{tabular}
\end{center}

\subsection{Device transfer}

\begin{center}
\begin{tabular}{@{}p{0.24\linewidth}p{0.70\linewidth}@{}}
\toprule
\textbf{Command} & \textbf{Purpose} \\
\midrule
\code{transfer copy} & Copy an active key identity between two devices
  that are both open. The source stays active. \\
\code{transfer move} & Move an active key identity between two open
  devices. The source becomes a ghost only after the destination
  commits. \\
\code{transfer recover} & Finish or roll back a transfer that stopped
  between the two devices. \\
\code{transfer enroll} & Record a slot on this device as an active key
  identity. \\
\code{transfer list} & List identities on a device (ghosts hidden unless
  \flag{--all}). \\
\code{transfer relay-send} & Seal a copy or move to the destination and
  leave the source prepared; the ghost applies only after the
  destination's acknowledgement is finalized. \\
\code{transfer relay-collect} & Pull a sealed copy or move, install it, and
  post the acknowledgement. \\
\code{transfer relay-finalize} & Apply a destination acknowledgement and
  finish the source. \\
\code{device init} & Create \file{device.kq} and an empty vault directory
  in this store's registry. \\
\code{device provision} & Add a passphrase-sealed slot to a container and
  print its public keys. \\
\code{device list} & List a container's slot labels and public keys. \\
\code{device bind} & Record a container's device id for a slot's
  registered keys. \\
\code{device register} & Register a slot's public key (a lone key file
  still uses top-level \code{register}). \\
\code{device publish} & Publish a container's public descriptor to the
  relay's device directory. Slot secrets stay local. \\
\code{device relay-relocate} & Seal one slot onto another device through
  the relay; prints the relocate id. The source slot stays until
  \code{relay-drop} sees the matching acknowledgement. \\
\code{device relay-accept} & Pull sealed relocates, install each slot, and
  post the acknowledgement. \\
\code{device relay-drop} & Delete the source slot after a verified
  relocate acknowledgement, keyed to the specific relocate id. \\
\bottomrule
\end{tabular}
\end{center}

\subsection{\code{keyquorum-device}: the container tool}

A separate binary for one physical container of logical identity slots,
used without an organization database (see
\hyperref[sec:devices-and-custody]{Devices and custody}). Private keys
stay in the slot token; this tool never prints them.

\begin{center}
\begin{tabular}{@{}p{0.24\linewidth}p{0.70\linewidth}@{}}
\toprule
\textbf{Command} & \textbf{Purpose} \\
\midrule
\code{init} & Create \file{device.kq} and an empty vault directory. \\
\code{provision} & Add a slot; prints its public keys and prompts for a
  passphrase. \\
\code{list} & List slot labels and public keys from \file{device.kq}. \\
\code{public} & Print one slot's public keys. \\
\code{unwrap-share} & Unwrap a sealed share with a slot's encryption key;
  the raw share prints to stdout. \\
\code{sign} & Sign stdin with a slot's signing key; the hex signature
  prints to stdout. \\
\code{relocate} & Move a slot onto another container. The keypair is
  unchanged; only the device id it is bound to changes. \\
\bottomrule
\end{tabular}
\end{center}
